\notechapter{N-20220806}{Set Theory Exercises, Closures, Functions, and Algebraic Side Notes}

\section{Images, preimages, and finite set examples}

The first page contains set-theoretic exercises.  For example, let
\[
  A=\{2,0,1,3\}
\]
and define
\[
  B=\{x:x\in A,\ 2-x^2\notin A\}.
\]
Checking each element:
\[
  x=-2\Rightarrow 2-x^2=-2\notin A,
\]
\[
  x=0\Rightarrow 2-x^2=2\in A,
\]
\[
  x=-1\Rightarrow 2-x^2=1\in A,
\]
\[
  x=-3\Rightarrow 2-x^2=-7\notin A.
\]
The page concludes with a finite set \(B\) after direct inspection.  The point is that a set defined by a condition is not mysterious: test the condition element by element when the universe is small.

\section{Closure under addition}

The next exercise asks for a subset \(S\subseteq\mathbb R\) closed under addition:
\[
  x,y\in S\quad\Rightarrow\quad x+y\in S.
\]
Examples include:
\[
  \mathbb Z,\qquad \mathbb Q,\qquad \mathbb R,\qquad [0,\infty).
\]
A finite example such as \(\mathbb Z/6\mathbb Z\) also works if addition is understood modulo \(6\).  The note correctly distinguishes ordinary addition from addition modulo an equivalence relation.

A related problem says: if \(S_1,S_2\subseteq\mathbb R\) are closed under addition and are proper in a suitable sense, prove there is an element \(c\) outside their union.  This resembles the Dedekind-cut intuition: two proper additive-closed pieces cannot cover the entire real line unless one of them already absorbs too much structure.

\section{A combinatorial maximal-set problem}

Another problem considers
\[
  M\subseteq\{1,2,\ldots,2011\}
\]
with the property that for any two elements \(a,b\in M\), at least one of
\[
  a\mid b,\qquad b\mid a,\qquad |a-b|=\max M
\]
holds.  The note first explores powers of \(2\):
\[
  1,2,2^2,\ldots,2^{10},
\]
since divisibility is automatic along a chain.  It then observes that if one tries to add a number not on the same divisibility chain, the condition quickly fails.

The key structural idea is that divisibility partially orders the integers.  A set in which every pair is comparable under divisibility is a chain.  Powers of a fixed prime give the cleanest example of a long chain.  This is an early encounter with posets and extremal combinatorics.

\section{Functional and algebraic exercises}

The later pages mix several algebraic problems.  One recurring theme is solving equations by rewriting them into factored or symmetric forms.  For instance, equations involving
\[
  x^2+y^2,\qquad xy,\qquad x+y
\]
are usually better handled through substitutions like
\[
  u=x+y,\qquad v=xy.
\]
This is the elementary symmetric-polynomial viewpoint.

Another page discusses functions and constraints of the form
\[
  f(A\cap B),\quad f(A\cup B),\quad f^{-1}(C).
\]
The essential warning is that images and preimages behave differently:
\[
  f^{-1}(C\cap D)=f^{-1}(C)\cap f^{-1}(D)
\]
always, while
\[
  f(A\cap B)=f(A)\cap f(B)
\]
need not hold unless \(f\) is injective on the relevant sets.

\section{Matrix and coordinate side computations}

Several pages contain coordinate computations and small matrix reductions.  These are not central enough to deserve separate chapters, but they reinforce a common habit: reduce a system to normal form.  Whether one is checking membership in a set, solving a function equation, or reducing a matrix, the practical goal is to isolate independent parameters.


\section{Images and preimages behave differently}

For a function \(f:X\to Y\), the image of a set \(A\subseteq X\) is
\[
  f(A)=\{f(a):a\in A\},
\]
while the preimage of \(B\subseteq Y\) is
\[
  f^{-1}(B)=\{x\in X:f(x)\in B\}.
\]
Preimages preserve unions, intersections, and complements exactly.  Images preserve unions but not intersections in general.  This asymmetry is why topology defines continuity using preimages of open sets.

\section{Closure operations}

Many set exercises secretly ask whether an operation is a closure operator.  A closure operator should be extensive, monotone, and idempotent:
\[
  A\subseteq \overline A,\qquad
  A\subseteq B\Rightarrow \overline A\subseteq \overline B,\qquad
  \overline{\overline A}=\overline A.
\]
Recognizing these three properties helps organize otherwise scattered set-manipulation problems.

\section{The reusable pattern}

This note is scattered, but its center is closure.  A set may be closed under addition, a family of elements may be closed under a divisibility relation, a function preimage is closed under Boolean operations, and a solution set may be closed under linear combinations.  Recognizing closure is one of the first steps toward algebraic structure.



\paragraph{Expanded synthesis.}
For this chapter, the elementary material should be read as a study of what remains invariant when coordinates change. The earlier sections (`Images, preimages, and finite set examples`, `Closure under addition`, `A combinatorial maximal-set problem`, `Functional and algebraic exercises`, `Matrix and coordinate side computations`) compute with arrays, bases, pivots, determinants, or eigenvectors; the deeper object is a linear map together with the spaces it acts on. A matrix is a representation of that map after bases have been chosen. This is why formulas such as
\[
  A' = P^{-1}AP,\qquad \widetilde A = T^{-1}AS
\]
are not cosmetic: they distinguish an intrinsic morphism from a coordinate description.

The structural hierarchy is as follows. Kernels record what a map forgets, images record what it reaches, cokernels record obstructions to solvability, and quotients record information after an equivalence has been imposed. Determinants act on the top exterior power and therefore measure oriented volume; traces are invariant under cyclic rearrangement and become characters in representation theory; adjoints depend on an inner product and lead to self-adjoint spectral theory.

A useful way to return to the computations is to ask, for every formula, which invariant it is measuring. Row reduction measures rank and kernel; diagonalization measures decomposition into invariant subspaces; Gram--Schmidt measures orthogonal projection; positive definiteness measures ellipsoid geometry; tensor notation measures how components transform. The elementary methods are therefore the finite-dimensional laboratory in which duality, quotienting, functoriality, and spectral decomposition can be seen clearly.


\section{Closure operators}

The closure-style exercises in the original note can be organized by the notion of a closure operator.  On a partially ordered set $P$, a closure operator is a map
\[
  c:P\to P
\]
satisfying
\[
  x\le c(x),\qquad x\le y\Rightarrow c(x)\le c(y),\qquad c(c(x))=c(x).
\]
These are called extensivity, monotonicity, and idempotence.

Many familiar constructions have this form:
\[
\begin{aligned}
  A&\mapsto \overline A &&\text{in topology},\\
  S&\mapsto \operatorname{span}(S) &&\text{in linear algebra},\\
  S&\mapsto \langle S\rangle &&\text{in group theory}.
\end{aligned}
\]
The elementary operation ``generate the smallest object containing a given set'' is exactly closure.

\section{Closed sets as fixed points}

An element $x$ is closed if
\[
  c(x)=x.
\]
Thus closed sets are fixed points of the closure operator.  The family of closed elements is closed under arbitrary meets.  In set language, topologically closed sets are stable under arbitrary intersections; subspaces generated by vector sets are stable under intersections; subgroups generated by sets are stable under intersections.

This explains a standard construction:
\[
  c(S)=\bigcap\{C:\ S\subseteq C,\ C\text{ closed}\}.
\]
The closure of $S$ is the intersection of all closed objects containing it.

\section{Galois connections}

A Galois connection between posets $P$ and $Q$ is a pair of order-reversing or order-preserving maps that translate inequalities in one side into inequalities on the other.  A common form is
\[
  f(p)\le q\quad\Longleftrightarrow\quad p\le g(q).
\]
The composite $g\circ f$ is then a closure operator.

In linear algebra, for a subset $S$ of vectors, one can define the annihilator
\[
  S^\perp=\{\varphi\in V^*:\ \varphi(s)=0\text{ for all }s\in S\}.
\]
Taking annihilators reverses inclusion and creates closure-like double-perp operations.

Thus the elementary habit of taking complements, generated sets, or common solution sets points toward a deep pattern: objects can be studied through the constraints they impose, and closure is often a double-constraint operation.

\section{Tarski fixed-point theorem}

On a complete lattice, every monotone map has a fixed point.  This is Tarski's fixed-point theorem.  A complete lattice is a partially ordered set in which every subset has a supremum and an infimum.

For a monotone map $F:L\to L$, the set of fixed points
\[
  \{x\in L:\ F(x)=x\}
\]
is itself a complete lattice.  This theorem appears in logic, semantics of programming languages, and inductive definitions.

The elementary idea of repeatedly applying a closure process until nothing changes is therefore a finite visible version of a fixed-point theorem for ordered structures.

\section{Returning to the original exercises}

Whenever the original note asks for a smallest set satisfying some property, an invariant collection, or a repeated operation that stabilizes, the advanced structure is closure.  The same pattern appears in topology, algebra, logic, and theoretical computer science.  The elementary set manipulations are the computational entrance to closure operators and fixed-point mathematics.

% Second-pass deepening for waves 2-4: wave 3 A-C part 1
\section{Generation as an intersection construction}

Many elementary problems ask for the smallest set satisfying a property.  The general construction is:
\[
  \langle S\rangle=\bigcap\{C:S\subseteq C,\ C\text{ has the desired closure property}\}.
\]
This works because the intersection of closed objects is closed.  For subspaces, intersections of subspaces are subspaces.  For subgroups, intersections of subgroups are subgroups.  For topological closed sets, arbitrary intersections are closed.

This explains why ``generated by'' is such a stable phrase across mathematics.  The generated object is not guessed; it is forced by closure under operations and minimality under inclusion.

\section{Closure and logic of consequences}

A closure operator can also be interpreted as logical consequence.  If $S$ is a set of assumptions, $c(S)$ is the set of all consequences of $S$.  Extensivity says assumptions are consequences of themselves.  Monotonicity says more assumptions give more consequences.  Idempotence says taking consequences twice adds nothing new.

This logic viewpoint connects the original set exercises to algebraic logic and model theory.  In algebra, closure under operations generates structures; in logic, closure under inference generates theories.
